% Part: normal-modal-logic
% Chapter: completeness
% Section: completeness-K

\documentclass[../../../include/open-logic-section]{subfiles}

\begin{document}

\olfileid[pt]{nml}{com}{cmk}

\olsection{Determinação e completude para \Log{K}}

Estamos agora preparados para usar o modelo canônico a fim de
estabelecer a completude. A completude decorre do fato de que as
!!{formula}s verdadeiras no modelo canônico de~$\Sigma$ coincidem
exatamente com aquelas !!{derivable} em~$\Sigma$. Diz-se que os modelos com essa
propriedade \emph{determinam}~$\Sigma$.

\begin{defn}
  Um modelo $\mModel{M}$ \emph{determina} uma lógica modal
  normal~$\Sigma$ precisamente quando $\mSat{M}{!A}$ se e somente se
  $\Sigma \Proves !A$, para todas as !!{formula}s~$!A$.
\end{defn}

\begin{thm}[Determinação]\ollabel{thm:determination}
   $\mSat{M^\Sigma}{!A}$ se e somente se $\Sigma \Proves !A$.
\end{thm}

\begin{proof}
  Se $\mSat{M^\Sigma}{!A}$, então, para todo conjunto completo
  $\Sigma$-consistente $\Delta$, temos
  $\mSat{M^\Sigma}{!A}[\Delta]$. Logo, pelo lema da verdade, $!A \in
  \Delta$ para todo conjunto completo $\Sigma$-consistente $\Delta$;
  daí, por \olref[lin]{cor:provability-characterization} (com $\Gamma =
  \emptyset$), $\Sigma \Proves !A$.

  Reciprocamente, se $\Sigma \Proves !A$, então, por
  \olref[ccs]{prop:ccs-properties}\olref[ccs]{prop:ccs-closed}, todo
  conjunto completo $\Sigma$-consistente $\Delta$ contém~$!A$ e,
  portanto, pelo lema da verdade, $\mSat{M^\Sigma}{!A}[\Delta]$ para
  todo~$\Delta \in W^\Sigma$, isto é, $\mSat{M^\Sigma}{!A}$.
\end{proof}

Como o modelo canônico de \Log{K} determina~\Log{K}, obtemos
imediatamente a completude de~\Log{K} como corolário:

\begin{cor}\ollabel{cor:Kcomplete}
  A lógica modal básica \Log{K} é completa com respeito à
  classe de todos os modelos, isto é, se $\Entails !A$, então $\Log{K}
  \Proves !A$.
\end{cor}

\begin{proof}
  Por contraposição, se $\Log{K} \Proves/ !A$, então, pela determinação,
  $\mSat/{M^{\Log{K}}}{!A}$ e, portanto, $!A$ não é válida.
\end{proof}

Para o caso geral da completude de um sistema $\Sigma$ com respeito a
uma classe de modelos, por exemplo, de \Log{KTB4} com respeito à classe
dos modelos reflexivos, simétricos e transitivos, a determinação por si
só não basta. Devemos mostrar também que o modelo canônico do sistema
$\Sigma$ é membro da classe, o que não decorre de maneira óbvia da
construção do modelo canônico---nem é sempre verdade!

\end{document}
